\documentclass[12pt]{article}
\usepackage[margin=1in]{geometry}
\usepackage{amsmath}
\usepackage{amssymb}
\usepackage{amsthm}
\usepackage{enumitem}
\usepackage{xcolor}
\usepackage{pagecolor}
\usepackage{marginnote}
\usepackage{fancyhdr}
\usepackage{bm}
\usepackage{etoolbox}
\usepackage{mdframed}

\definecolor{cream}{HTML}{faf7f1}
\pagecolor{cream}
\mdfdefinestyle{lemmastyle}{backgroundcolor=cream}

\newcommand{\defn}[1]{\textbf{#1}}
\newcommand{\defeq}{\mathrel{\vcenter{\baselineskip0.5ex \lineskiplimit0pt
                     \hbox{\scriptsize.}\hbox{\scriptsize.}}}=}
\newcommand{\T}{\top}
\newcommand{\F}{\bot}
\newcommand{\Pow}{\mathbb{P}}
\newcommand{\suc}{\mathfrak{suc}}
\newcommand{\N}{\mathbb{N}}
\newcommand{\Z}{\mathbb{Z}}
\newcommand{\Func}{\mathcal{F}}

\newtheoremstyle{proofstyle}
  {}{}{}{}{\bfseries}{.}{ }{}
\theoremstyle{proofstyle}
\newtheorem*{proof*}{Proof}

\newcommand{\sidenote}[1]{\marginnote{\small\itshape #1}}

\makeatletter
\renewcommand{\maketitle}{%
  \begin{flushleft}
    {\Large\bfseries CS173: Discrete Structures}\\[0.3em]
    {\Large\bfseries Problem Set 6}\\[0.5em]
    {\normalsize Mohammed Al Salhi}
  \end{flushleft}
  \vspace{1em}
}
\makeatother

\AtBeginDocument{\mathversion{bold}}

\begin{document}

\maketitle

\begin{enumerate}[left=0pt, itemsep=2em, parsep=1em]

%% ---------------------------------------------------------------
%% Problem 1
%% ---------------------------------------------------------------
\item Prove $(\forall a, b, x, y, n \in \Z)((n \mid x \land n \mid y) \Rightarrow (n \mid ax + by))$.

\begin{proof*}
Let $a, b, x, y, n \in \Z$ be arbitrary, and suppose $n \mid x$ and $n \mid y$.

By definition of divisibility, there exist integers $k_1, k_2 \in \Z$ such that $x = n k_1$ and $y = n k_2$. Then:
\begin{align*}
  ax + by &= a(nk_1) + b(nk_2) \\
          &= n(ak_1) + n(bk_2) \\
          &= n(ak_1 + bk_2)
\end{align*}
Since $\Z$ is closed under addition and multiplication, $ak_1 + bk_2 \in \Z$.

Therefore $n \mid ax + by$ by definition of divisibility. Since $a, b, x, y, n$ were arbitrary, we conclude $(\forall a, b, x, y, n \in \Z)((n \mid x \land n \mid y) \Rightarrow (n \mid ax + by))$.
\end{proof*}

%% ---------------------------------------------------------------
%% Problem 2
%% ---------------------------------------------------------------
\item We say an integer $z \in \Z$ is \defn{even} $:\Leftrightarrow 2 \mid z$, and we say $z$ is \defn{odd} $:\Leftrightarrow 2 \mid z - 1$.

\begin{enumerate}[label=(\alph*), itemsep=1.5em, leftmargin=2em]

  \item Prove that no integer can be both even and odd at the same time.

  \begin{proof*}
  Suppose for contradiction that some $z \in \Z$ is both even and odd.

  Then $2 \mid z$ and $2 \mid z - 1$, so there exist $k_1, k_2 \in \Z$ such that $z = 2k_1$ and $z - 1 = 2k_2$. Substituting the first into the second:
  \begin{align*}
    2k_1 - 1 &= 2k_2 \\
    2(k_1 - k_2) &= 1
  \end{align*}
  Since $\Z$ is closed under subtraction, $k_1 - k_2 \in \Z$, this gives $2 \mid 1$, which contradicts the fact that $1 < 2$. Therefore no integer can be both even and odd.
  \end{proof*}

  \item Prove that $(\forall z \in \Z)((z \text{ is even}) \Leftrightarrow (z + 1 \text{ is odd}))$.

  \begin{proof*}
  Let $z \in \Z$ be arbitrary.

  $(\Rightarrow)$ Suppose $z$ is even, i.e.\ $2 \mid z$. We must show $z + 1$ is odd, i.e.\ $2 \mid (z+1) - 1$. But $(z+1) - 1 = z$, and $2 \mid z$ by assumption. Therefore $z + 1$ is odd.

  $(\Leftarrow)$ Suppose $z + 1$ is odd, i.e.\ $2 \mid (z+1) - 1$. But $(z+1) - 1 = z$, so $2 \mid z$. Therefore $z$ is even.

  Since $z$ was arbitrary, $(\forall z \in \Z)((z \text{ is even}) \Leftrightarrow (z + 1 \text{ is odd}))$.
  \end{proof*}

  \item Prove that $(n \text{ is even}) \vee (n \text{ is odd})$ for all $n \in \N$.

  \begin{proof*}
  We proceed by weak induction on $n$.

  \textbf{Basis Step:} We must show $0$ is even or $0$ is odd. Since $0 = 2 \cdot 0$ and $0 \in \Z$, we have $2 \mid 0$, so $0$ is even.

  \textbf{Inductive Hypothesis:} Let $k \in \N$ and assume $k$ is even or $k$ is odd.

  \textbf{Inductive Step:} We must show $k+1$ is even or $k+1$ is odd. We consider two cases.

  \textit{Case 1:} Suppose $k$ is even. By part (b), $k + 1$ is odd.

  \textit{Case 2:} Suppose $k$ is odd, so $2 \mid k - 1$. Then there exists $m \in \Z$ such that
  \begin{align*}
        k - 1 &= 2m \\
        k &= 2m + 1 \\
        k + 1 &= 2m + 2 \\
        k + 1 &= 2(m+1)
  \end{align*}
  Since $m + 1 \in \Z$, we have $2 \mid k+1$, so $k+1$ is even.

  Therefore, by weak induction, $(n \text{ is even}) \vee (n \text{ is odd})$ for all $n \in \N$.
  \end{proof*}
\end{enumerate}

\newpage

%% ---------------------------------------------------------------
%% Problem 3
%% ---------------------------------------------------------------
\item Show that $(\exists m \in \N)(\forall n \in \N)(n \geq m \Rightarrow n^2 < 2^n)$.

\begin{proof*}
By inspection of small cases, $m = 5$ suffices. We claim $m = 5$ witnesses the existential.

We first establish a lemma.
\begin{mdframed}[style=lemmastyle]
\textbf{Lemma:} $(\forall k \in \N)(k \geq 5 \Rightarrow 2k + 1 < k^2)$.

\textit{Proof of Lemma.} We proceed by weak induction on $k$.

\textit{Basis Step:} We have $2(5) + 1 = 11 < 25 = 5^2$.

\textit{Inductive Hypothesis:} Let $k \in \N$ with $k \geq 5$, and assume $2k + 1 < k^2$.

\textit{Inductive Step:} We must show $2(k+1) + 1 < (k+1)^2$. Observe:
\begin{align*}
  2(k+1) + 1 &= 2k + 3 \\
  (k+1)^2    &= k^2 + 2k + 1
\end{align*}
So it suffices to show $2k + 3 < k^2 + 2k + 1$, i.e.\ $2 < k^2$. Since $k \geq 5$, we have $k^2 \geq 25 > 2$. Therefore $2(k+1) + 1 < (k+1)^2$. \qed
\end{mdframed}

We proceed by weak induction on $n$.

\textbf{Basis Step:} We have $5^2 = 25 < 32 = 2^5$.

\textbf{Inductive Hypothesis:} Let $k \in \N$ with $k \geq 5$, and assume $k^2 < 2^k$.

\textbf{Inductive Step:} We must show $(k+1)^2 < 2^{k+1}$. Observe:
\begin{align*}
  (k+1)^2 &= k^2 + 2k + 1 \\
           &< 2k^2           &&\text{by the Lemma, since }2k+1 < k^2 \\
           &< 2 \cdot 2^k &&\text{by the inductive hypothesis} \\
           &= 2^{k+1}        &&\text{by definition of exponentiation}
\end{align*}
Therefore $(k+1)^2 < 2^{k+1}$.

By induction, $(\forall n \in \N)(n \geq 5 \Rightarrow n^2 < 2^n)$, and so $m = 5$ witnesses $(\exists m \in \N)(\forall n \in \N)(n \geq m \Rightarrow n^2 < 2^n)$.
\end{proof*}

\newpage

%% ---------------------------------------------------------------
%% Problem 4
%% ---------------------------------------------------------------
\item The factorial function ${\cdot}! : \N \to \N$ is the recursive function given below for all $n \in \N$.
\[
  0! \defeq 1
\]
\[
  (n+1)! \defeq (n!) \cdot (n+1)
\]
Prove that $(\forall n \in \N)(n \geq 4 \Rightarrow 2^n < n!)$.

\begin{proof*}
We will prove by weak induction on $n$.

\textbf{Basis Step:} We must show $2^4 < 4!$. We first expand $4!$ using the recursive definition.
\begin{align*}
  4! &= ((((1) \cdot 1) \cdot 2) \cdot 3) \cdot 4 &&\text{by definition of }! \\
     &= 24 &&\text{by arithmetic}
\end{align*}
Since $2^4 = 16 < 24 = 4!$, the basis step holds.

\textbf{Inductive Hypothesis:} Let $k \in \N$ with $k \geq 4$, and assume $2^k < k!$.

\textbf{Inductive Step:} We must show $2^{k+1} < (k+1)!$. Observe:
\begin{align*}
  2^{k+1} &= 2^k \cdot 2 &&\text{by definition of exponentiation} \\
          &< k! \cdot 2 &&\text{by the inductive hypothesis} \\
          &< k! \cdot (k+1) &&\text{since }k \geq 4 \text{ implies } k + 1 > 2\\
          &= (k+1)! &&\text{by definition of }!
\end{align*}
Therefore $2^{k+1} < (k+1)!$.

By induction, $(\forall n \in \N)(n \geq 4 \Rightarrow 2^n < n!)$.
\end{proof*}

\newpage

%% ---------------------------------------------------------------
%% Problem 5
%% ---------------------------------------------------------------
\item The Fibonacci sequence is the function $\mathcal{F} : \N \to \N$ defined below for all $n \in \N$.
\[
  \mathcal{F}(0) \defeq 0
\]
\[
  \mathcal{F}(1) \defeq 1
\]
\[
  \mathcal{F}(n+2) \defeq \mathcal{F}(n+1) + \mathcal{F}(n)
\]
Prove that $\displaystyle\sum_{i=1}^{n} \mathcal{F}(2i-1) = \mathcal{F}(2n)$ for all $n \in \N$.

\begin{proof*}
We proceed by weak induction on $n$.

\textbf{Basis Step:} We must show $\displaystyle\sum_{i=1}^{0} \mathcal{F}(2i-1) = \mathcal{F}(0)$. The left side is an empty sum, which equals $0$ by convention. The right side is $\mathcal{F}(0) = 0$ by definition. Therefore both sides are equal.

\textbf{Inductive Hypothesis:} Let $k \in \N$ and assume $\displaystyle\sum_{i=1}^{k} \mathcal{F}(2i-1) = \mathcal{F}(2k)$.

\textbf{Inductive Step:} We must show $\displaystyle\sum_{i=1}^{k+1} \mathcal{F}(2i-1) = \mathcal{F}(2k+2)$. Observe:
\begin{align*}
  \sum_{i=1}^{k+1} \mathcal{F}(2i-1)
    &= \left(\sum_{i=1}^{k} \mathcal{F}(2i-1)\right) + \mathcal{F}(2(k+1)-1)
      &&\text{by definition of iterated summation} \\
    &= \mathcal{F}(2k) + \mathcal{F}(2(k+1)-1)
      &&\text{by the inductive hypothesis} \\
    &= \mathcal{F}(2k) + \mathcal{F}(2k+1)
      &&\text{by arithmetic} \\
    &= \mathcal{F}(2k+2)
      &&\text{by definition of }\mathcal{F}
\end{align*}
Therefore $\displaystyle\sum_{i=1}^{k+1} \mathcal{F}(2i-1) = \mathcal{F}(2k+2)$.

By induction, $\displaystyle\sum_{i=1}^{n} \mathcal{F}(2i-1) = \mathcal{F}(2n)$ for all $n \in \N$.
\end{proof*}

\end{enumerate}

\end{document}